
\documentclass{article} % For LaTeX2e
\usepackage{iclr2027_conference,times}

% Optional math commands from https://github.com/goodfeli/dlbook_notation.
\input{math_commands.tex}

\usepackage{hyperref}
\usepackage{url}
\usepackage{booktabs}
\usepackage{amssymb}
\usepackage{graphicx}
\usepackage{enumitem}

\newcommand{\ie}{i.e.}\xspace
\newcommand{\eg}{e.g.}\xspace

\definecolor{pink}{RGB}{204, 0, 102}
\newcommand{\AR}[1]{\textcolor{pink}{\textit{(Angie: #1)}}}


% \title{Signal-Aware Framework for Multilingual Language Model Evaluation}
\title{Decision Reliability in Multilingual Model Development}

% Alternative titles:
% Predicting Reliable Decisions in Multilingual Model Scaling
% Decision Reliability in Multilingual Model Development/Scaling

% Evaluating Multilingual Evaluations: In Search for Benchmarks that Enable Reliable Decisions

% MAIN RESEARCH QUESTION
% People use evaluations of cheaper, small/partially trained models to make expensive training decisions. When is that practice actually reliable in multilingual model development?



% Authors must not appear in the submitted version. They should be hidden
% as long as the \iclrfinalcopy macro remains commented out below.
% Non-anonymous submissions will be rejected without review.

\author{María Grandury
EPFL\\
\And
Angelika Romanou
EPFL\\
\And
Clara Meister
EPFL
\And
Éléonore Hasler
\And
Antoine Bosselut
}

% The \author macro works with any number of authors. There are two commands
% used to separate the names and addresses of multiple authors: \And and \AND.
%
% Using \And between authors leaves it to \LaTeX{} to determine where to break
% the lines. Using \AND forces a linebreak at that point. So, if \LaTeX{}
% puts 3 of 4 authors names on the first line, and the last on the second
% line, try using \AND instead of \And before the third author name.

\newcommand{\fix}{\marginpar{FIX}}
\newcommand{\new}{\marginpar{NEW}}
\newcommand{\update}{\marginpar{UPDATE}}
\newcommand{\citeneeded}[1]{\textcolor{red}{[CITE: #1]}}
\newcommand{\checknum}[1]{\textcolor{orange}{#1}}

% TODO: Update constants
\newcommand{\nbenchmarks}{14 }
\newcommand{\nbenchmarktasks}{XX }
\newcommand{\nmodels}{162 }
\newcommand{\nckpts}{3000 }


%\iclrfinalcopy % Uncomment for camera-ready version, but NOT for submission.
\begin{document}
\maketitle
\begin{abstract}
Training multilingual language models requires repeated decisions about data mixtures and architectures, typically guided by evaluations of small proxy models.
Such evaluations are useful, however, only if results from small-scale models reliably follow the same trends as their larger-scale counterparts. 
While the conditions under which small-scale evaluations are predictive have been studied extensively in monolingual settings, these findings do not necessarily transfer to multilingual training, where model capacity is shared across languages and scaling behavior can depend strongly on the composition of the training mixture. 
%a challenge compounded in multilingual settings where many language–specific tasks exhibit little measurable signal at small scales.
This work asks four foundational questions around this practice: (1) for multilingual language models, at what model size do evaluations become sufficiently predictive of larger models? (2) how early during pretraining can decisions be taken? (3) can we determine before doing many expensive evaluations whether a particular evaluation is likely to be informative at small-scale? and (4) how do these patterns change as multilingual training data varies in composition? 
We pretrain a controlled ladder of \nmodels models spanning six sizes from 90M to 1.7B non-embedding parameters and six language settings from 1 to 50 languages, with four design interventions and seed replicates. 
We evaluate 12 aligned checkpoints from each run using per-language bits per byte over 100 languages and tasks from \nbenchmarks multilingual benchmark families.
% RQ1
We characterize which benchmark families exhibit predictable scaling across model size and training, and distinguish benchmarks whose performance improves consistently from those whose scaling is weak, noisy, or non-monotonic. 
% RQ2
For each benchmark, we quantify how well performance at smaller model sizes predicts performance at larger sizes, and how early in training this correspondence emerges. We then test whether these small-scale evaluations recover the relative effects of our data-mixture and architecture interventions observed at full scale.
% RQ3
Finally, we ask whether the reliability of a small-scale evaluation can itself be anticipated cheaply, comparing measurements that assess whether a benchmark has sufficient signal to recover full-scale results across model scales and as a function of the number and composition of training languages.
% RQ4
% Finally, we test whether scaling behavior learned from observed languages transfers to held-out languages, and how this transfer depends on the underlying design intervention. 
% We find that likelihood-based measurements scale predictably while four-option knowledge benchmarks sit at chance up to 1.7B; that a 350M proxy at 20\% of its run already reads a data decision the way the 1.7B reference does on bits per byte, while the benchmark suite reads every decision as a coin flip; that the reliability of a proxy is set by the size of the intervention's effect relative to seed noise rather than by the intervention itself; and that scaling exponents transfer across languages, so that one small measurement predicts a never-trained language's bits per byte at 1.7B within 7\%. 
We release the ladder and the per-checkpoint measurements to support cheaper, better-informed multilingual model development.
\end{abstract}

\section{Introduction}
\label{sec:introduction}

Pretraining multilingual language models requires billions of tokens of data and hundreds of thousands of euros in compute.
Modern post-training projects iterate over dozens of recipes
%for SFT mixes, preference data, and reinforcement learning settings 
\citep{lambert_tulu_2025}, each evaluated on long, expensive benchmark suites. Benchmarks inform decisions that impact the final model quality, but they vary in their diagnostic value. Some are sensitive to recipe changes, others have high variance, redundancy, or weak correlation with downstream objectives.
Improvements on certain benchmarks, therefore, may not reflect meaningful progress.
Evaluation cost scales with the number of decisions made, and unreliable benchmarks waste a large share of the development budget.
In particular, long-context extension produces effective windows shorter than the advertised ones
\citep{hsieh_ruler_2024}, and the later post-training stages, namely
supervised fine-tuning (SFT), direct preference optimization (DPO), and Reinforcement Learning with Verifiable Rewards (RLVR),
often score open-ended tasks with biased LLM judges \citep{dubois_length-controlled_2024, feuer_style_2025}.
%, which suffer from biases such as length preference \cite{dubois2025lengthcontrolledalpacaevalsimpleway} and style over substance \cite{feuer2025styleoutweighssubstancefailure}. The problem is more severe in multilingual settings, where judge competence drops on lower-resource languages \cite{hada2024largelanguagemodelbasedevaluators, watts2024parikshalargescaleinvestigationhumanllm}.
%Practitioners therefore face a tradeoff between evaluation coverage and efficiency without clear guidance on which benchmarks provide a reliable signal.
Moreover, these reliability gaps widen for multilingual settings \citep{kim_one_2025, hada_are_2024, watts_pariksha_2024}.

% As a result, practitioners often evaluate on extensive benchmark suites with diminishing diagnostic returns, or rely on small subsets of tasks whose signals may be noisy or misaligned with the decisions being made. Systematically identifying benchmarks that provide meaningful, stage-specific diagnostic signal could substantially improve the efficiency and reliability of multilingual model development.


Recent work has begun to examine the diagnostic value and efficiency of language model benchmarks, recognizing that not all benchmarks provide equally informative signals for guiding model development. \citet{heineman_signal_2025} introduce a signal-to-noise ratio (SNR) framework that quantifies benchmark reliability for English models. 
%In this framework, signal measures how well a benchmark separates models of similar scale, while noise captures variability in scores across training checkpoints or stochastic runs. 
Their empirical analysis shows that benchmarks with higher SNR more consistently preserve model rankings across scales, improving the accuracy of training decisions and reducing error when extrapolating performance using scaling laws. This work highlights that benchmark choice directly affects the reliability of decisions made during model development.
Other work targets evaluation efficiency, identifying benchmark subsets with high discriminative power \citep{chen_scaling_2025, gupta_improving_2025, zhou_lost_2026}.
%\citet{chen2025scalinglawspredictingdownstream} show that only some downstream tasks follow predictable scaling trends. \citet{gupta-etal-2025-improving-smart} propose SMART filtering to remove redundant, contaminated, or overly easy examples. \citet{Zhou_Huang_Zhao_Han_Wang_Chen_Yang_Bao_Dong_Xu_Zhu_Cao_Zhao_2026_lostinbenchmarks} use item response theory to identify items with weak discriminative power.

Together, these works suggest that benchmark effectiveness depends not only on coverage but also on how much reliable signal a benchmark provides relative to its evaluation cost. However, existing studies focus on English benchmarks and pretraining. 
It is unclear whether their conclusions transfer to multilingual model
development. Decision accuracy may differ across languages, since languages
differ in data availability, tokenization cost, and scaling behavior
\citep{petrov_language_2023, ahia_all_2023}. Cross-lingual transfer further
couples decisions across languages, since adding data for one language can help or hurt others depending on model capacity and language similarity \citep{chang_when_2023, longpre_atlas_2026}, so a benchmark that is reliable for one language may be uninformative for another.

%In this report, we present our progress in extending the aforementioned frameworks to multilingual settings to identify the most reliable multilingual benchmarks for each training stage that correlate with final model performance. This will enable practitioners to make better, cheaper decisions during multilingual model development.

%\begin{itemize}
%    \item Which multilingual benchmarks produce reliable model rankings at small scale, such that those rankings can inform model design decisions during early-stage ablations and generalize to larger models?
%    \item Can benchmark metrics, computed at small scale, serve as a proxy for decision accuracy, removing the need to evaluate large models?
%    \item Which benchmark design choices, including data source and curation, lead to higher decision accuracy? %Types of questions (driving license, topics), hard questions, filtering decisions, etc -> what makes a good question
%\end{itemize}

In this report, we extend the SNR framework to multilingual settings and ask
three questions, namely which multilingual benchmarks produce small-scale
rankings that hold at large scale, whether SNR computed at small scale can
serve as a proxy for decision accuracy, and which benchmark design choices
drive reliability. Using 36 custom pretrained models and 68 open-source
models from 270M to 70B parameters, we find that most multilingual
benchmarks are unmeasurable at small scale, that the most transferable
benchmark differs by language, with human-translated commonsense tasks
leading and Swahili left effectively unmeasured, and that SNR predicts
decision accuracy across two very different model populations. We release a
per-language reliability map that practitioners can consult for the
languages they target. We aim to enable practitioners to make better, cheaper decisions during multilingual model development.

\section{Related Work}

\subsection{Benchmark Reliability}

Recent work shows that benchmarks differ substantially in how reliably they guide language model development. \citet{heineman_signal_2025} introduce a signal-to-noise ratio (SNR) framework where signal measures a benchmark's ability to separate models and noise captures score variability across checkpoints or stochastic runs. They show that high-SNR benchmarks better preserve rankings and reduce scaling-law prediction error. Complementary work studies benchmark quality through ranking consistency, discriminability, and capability-alignment deviations \citep{qian_benchmark2_2026}, while IRT-based methods show that many benchmark items have weak discriminative power and can often be replaced by smaller, more informative subsets \citep{zhou_lost_2026,li_adaptive_2025}.

Efficiency-oriented work reaches similar conclusions.
\citet{perlitz_efficient_2024} quantify the reliability cost of benchmark design decisions and show that large suites can be cut by two orders of magnitude with little loss in ranking reliability. \citet{polo_tinybenchmarks_2024} use item response theory to select small representative subsets, estimating MMLU performance from 100 curated examples.
\citet{gupta_improving_2025} propose SMART filtering to remove easy, contaminated, or redundant examples, reducing evaluation size while preserving agreement with independent model rankings. Scaling-law studies further suggest that only some evaluations provide predictable development signal: \citet{chen_scaling_2025} forecast downstream performance from smaller models, while \citet{schaeffer_pretraining_2026} study scaling laws for generative metrics such as pass@$k$. Together, these papers motivate our focus on identifying which multilingual measurements remain reliable across model sizes and nearby pretraining checkpoints, rather than treating all benchmark averages as equally meaningful.

\subsection{Multilingual Evaluation}

Multilingual benchmarks inherit reliability problems from their English counterparts and add new ones. Most popular benchmarks are machine-translated from English, which introduces translation artifacts, cultural mismatch, and label noise that distort cross-lingual rankings \citep{plaza_spanish_2024, singh_global_2024, romanou_include_2024}. Tokenizer performance varies by an order of magnitude across languages, inflating both training and evaluation cost for non-Latin scripts and morphologically rich languages \citep{petrov_language_2023, ahia_all_2023}. Cross-lingual transfer further couples decisions across languages, since adding data for one language can help or hurt others depending on model capacity and language similarity \citep{chang_when_2023, longpre_atlas_2026}. These factors suggest that English reliability estimates do not transfer directly across languages, and motivate a multilingual-specific framework built on controlled pretraining runs.

\section{Methodology}
\label{sec:methodology}

We extend the Signal and Noise framework of \citet{heineman_signal_2025} from English to multilingual language model development. We first define decision accuracy (DA) as whether benchmark rankings transfer from smaller to larger models. We propose the signal-to-noise ratio (SNR) as a proxy for decision accuracy, and present multiple possible definitions for signal and noise. We then enumerate the models used to compute the correlations between DA and SNR, and the benchmark suite.

\subsection{Decision Accuracy}
\label{subsec:methodology-da}

Decision accuracy (DA) measures whether the ranking induced by small proxy models matches the ranking of larger target models on a benchmark. We define it as
\begin{equation}
\text{DA} = \frac{1}{|\mathcal{P}|} \sum_{(a,b) \in \mathcal{P}} \mathbb{1}\big[\text{sign}(B(s_a) - B(s_b)) = \text{sign}(B(m_a) - B(m_b))\big]
\end{equation}
where $\mathcal{P}$ is a set of model pairs, $s_a$ is the small proxy model trained with characteristic $a$, $m_a$ is the corresponding target model, and $B(\cdot)$ returns the benchmark score. A higher DA means rankings observed at small scale transfer reliably to large scale, and therefore that the benchmark can guide development decisions.

\textbf{Pretraining DA.} For pretraining, we propose two variants. DA by size compares models from the same family across two sizes, for example 1B and 8B parameters. DA by checkpoint compares intermediate checkpoints of a single model, taken at 10\%, 33\%, 66\%, or 100\% of training in terms of FLOPS.

%\textbf{Post-training DA.} For post-training, we also propose two variants. Cross-context DA measures whether short-context scores preserve recipe rankings after context extension. Cross-stage DA measures whether checkpoint scores preserve rankings across the SFT, DPO, and RLVR stages.

\textbf{From DA to SNR.} Decision accuracy is expensive to compute, since it requires evaluating the larger target models. We therefore search for cheaper benchmark statistics whose values correlate with DA, computed at the benchmark level and at the sample level, for each language and each stage. The next subsection introduces the candidates.

\subsection{Signal and Noise}
\label{subsec:methodology-snr}

The signal-to-noise ratio (SNR) summarizes a benchmark's reliability as the ratio between signal, how well a benchmark separates models of similar compute, and noise, the variability in scores across training checkpoints or stochastic runs. 
% for posttraining: signal, the variance of mean scores across recipes, and noise, the variance across checkpoints and seeds for a fixed recipe.
We define it as
\begin{equation}
\text{SNR} = \frac{\text{Signal}}{\text{Noise}}.
\end{equation}
A higher SNR means score differences across models are large relative to the benchmark's intrinsic variability. The framework is not tied to a particular signal or noise definition, so we evaluate several candidates and select the pair with the strongest empirical correlation to DA.

\textbf{Signal candidates.} Given the final checkpoints of training runs at similar compute $C_\text{final}$, we consider five candidate signal definitions to capture a different notion of how spread out model scores are. Variance is the average squared distance from the mean. Mean distance is the average pairwise distance between points. Relative standard deviation is the standard deviation divided by the mean. Star discrepancy is the largest difference between any point and the uniform distribution. Dispersion is the largest pairwise distance,
\begin{equation}
\text{Dispersion}(C_\text{final}) = \max_{i \neq j} \|c_i - c_j\|.
\end{equation}

\textbf{Noise candidates.} We consider four sources of noise that the framework can in principle account for. Seed noise is the the standard deviation of the final checkpoint across training runs with different random seeds. Data order noise is the equivalent across runs that vary the order of training documents. Total variation is the average change in metric score across training checkpoints minus an improvement term. Checkpoint-to-checkpoint noise is proposed by \citet{heineman_signal_2025} as a proxy when the first three are too expensive to estimate at large scale and depends only on the final $n$ training checkpoints:
\begin{equation}
\text{Checkpoint noise} = \sigma\big(\{U(t_j)\}_{j=T-k+1}^{T}\big).
\end{equation}
%All four require access to multiple training runs or intermediate checkpoints, which limits their applicability to external models.
We propose a fifth noise definition, computable from a single evaluation run on any public model. Benchmark noise is the average relative standard deviation of scores across $k = 5$ disjoint folds of the evaluation set:
\begin{equation}
\text{Noise}(M) = \frac{1}{|M|} \sum_{m \in M} \frac{\sqrt{\tfrac{1}{k}\sum_{i=1}^{k}(m_i - \bar{m})^2}}{\bar{m}}
\end{equation}
where $m_1, \dots, m_k$ are the per-fold scores of model $m$. The fold scores are obtained by partitioning the stored per-question outputs of a single evaluation run, so no inference is repeated and no intermediate checkpoints are needed. This makes the framework applicable to any model with public scores.


\subsection{Models}
\label{subsec:methodology-models}

We pretrain XX small models across six scales, 90M, 175M, 350M, 600M, 1B, and 1.7B non-embedding parameters, using XX English/multilingual data mixtures ... We extend the analysis to 8B and 70B models by considering open-source models with available intermediate checkpoints.

\textbf{Custom proxy models.} 
All models follow a decoder-only transformer architecture with grouped query attention, rotary positional embeddings (RoPE) \citep{su_roformer_2023} and XiELU as the activation function \citep{huang_xielu_2025}. The architecture scales by jointly increasing depth, model dimension, and number of attention heads while keeping the KV head ratio fixed at 1/4. All models are trained on approximately 100B tokens drawn from FineEdu-DCLM and FineWeb2 \citep{penedo_fineweb_2024}, with a shared vocabulary of 131,072 tokens. Since this vocabulary size places a large fraction of parameters in the embeddings at these scales, we tie input and output embeddings. Training is done in BF16 mixed precision using Megatron-LM \citep{shoeybi_megatron-lm_2019}.


\textbf{Data mixtures.} We construct 3 mixtures by varying the ratio of English to multilingual data, namely 90/10, 60/40, and 30/70. The English component comes from FineWeb-Edu \citep{lozhkov_fineweb-edu_2024} filtered through the DCLM pipeline \citep{li_datacomp-lm_2025}. The multilingual component is drawn from FineWeb2 \citep{penedo_fineweb2_2025} using the Apertus high-quality filter \citep{messmer_enhancing_2026}, covering the top 200 languages and preserving the original language distribution of the corpus.

\textbf{Open-source models.} We extend the analysis to larger scales and later training stages with three open-source multilingual model families (Apertus \citep{apertus_apertus_2025}, SmolLM3 \citep{bakouch_smollm3_2025}, and OLMo3 \citep{olmo_olmo_2025}) with intermediate checkpoints across post-training stages and sizes from 0.6B to 8B.
%For pretraining, we add Apertus 1B, 3B, and 8B base. For mid-training, we include SmolLM3 3B Base, OLMo3 7B Base, and Apertus 8B and 70B base. For post-training, Apertus 0.6B and 1.7B distilled and distilled SFT, SmolLM3 3B Instruct, OLMo3 7B SFT, DPO, and Instruct (RLVR), and Apertus 8B and 70B Instruct.


\subsection{Benchmarks}
\label{subsec:methodology-benchmarks}

\textbf{Languages.} We evaluate on a selection of 12 languages that balances typological breadth with benchmark availability (Table \ref{tab:languages}).
%Following the resource taxonomy of \citet{joshi-etal-2020-language-resources},
We include languages from across the resource spectrum, %, from high -resource (English, Spanish, Russian, Chinese, Japanese, Arabic) through mid-resource (Hindi, Turkish, Vietnamese) to lower-resource (Swahili, Basque).
covering seven language families and one language isolate \citep{campbell_language_2010},
%: Indo-European with four branches (Germanic, Romance, Slavic, Indo-Aryan), Sino-Tibetan, Japonic, Afro-Asiatic, Turkic, Niger-Congo, Austroasiatic, and Basque.
six writing systems, % are represented
%: Latin, Cyrillic, Han, Devanagari, Arabic, and Japanese mixed scripts
%The set also covers
and typological features that stress tokenization and modeling \citep{petrov_language_2023, ahia_all_2023, limisiewicz_tokenization_2023}.
%\citet{petrov2023tokenizersunfairness} and \citet{ahia-etal-2023-tokenization-cost} show that the same content can require many more tokens in some languages than in others.
%, with differences of up to 15 times across the languages of the FLORES-200 benchmark CITE.
%This gap is largest for non-Latin scripts and for morphologically rich languages, and it directly affects training compute, context length, and inference cost. \citet{limisiewicz-etal-2023-tokenization-vocabulary} further show that vocabulary allocation and overlap across languages systematically influence downstream task performance. Our set covers the main axes along which these effects vary: agglutinative morphology (Turkish, Basque), absence of explicit word boundaries (Chinese, Japanese), root-and-pattern morphology (Arabic), tonal systems (Chinese, Vietnamese), and noun-class agreement (Swahili). Basque, as a language isolate \citet{Campbell2010LanguageIsolates}, gives a clean diagnostic case in which transfer from typologically related high-resource languages is not available.


%We intentionally select benchmarks that are widely adopted in the evaluation of each training stage.
\textbf{Benchmark suite.} We evaluate pretraining benchmarks covering several capabilities:
world knowledge, reading comprehension, natural language inference, commonsense reasoning (decomposed into physical, narrative, and activity-based variants), grammatical acceptability, and truthfulness \citep{conneau_xnli_2018, bandarkar_belebele_2024, zellers_hellaswag_2019}. 
%World knowledge is measured with Global MMLU \citep{singh2024globalmmluunderstandingaddressing} and the multilingual MMLU and ARC of the Okapi suite \citep{}; reading comprehension with Belebele \citep{bandarkar-etal-2024-belebele}; and natural language inference with XNLI \citep{conneau-etal-2018-xnli}. For commonsense reasoning we use three complementary benchmarks: Global PIQA for physical commonsense, XStoryCloze for narrative commonsense \citep{}, and the multilingual HellaSwag for activity-based commonsense CITE. Truthfulness is measured with the multilingual TruthfulQA \citep{}, and mathematical reasoning with MGSM in its direct (no chain-of-thought) variant CITE, the form usable on pretrained models that do not yet produce reliable chain-of-thought outputs. Grammatical acceptability is measured with MultiBLiMP \citep{}. The combination of capabilities is intentional: grammatical acceptability emerges early in pretraining and provides signal at small model scales, while knowledge and mathematical reasoning emerge later and discriminate larger models.
% - add egyhellaswag?
% - turkishmmlu?
The mid-training suite reuses some pretraining benchmarks and focuses on regional knowledge, reasoning capabilities, including math, code, and long-context tasks \citep{singh_global_2024,romanou_include_2024,shi_language_2022}.
The post-training suite considers reasoning, code, and adds instruction-following, bias, and safety benchmarks \citep{zhang_p-mmeval_2025, dussolle_m-ifeval_2025, nangia_crows-pairs_2020}. Table \ref{tab:benchmarks} shows the language coverage of the benchmarks.
%The multilingual benchmark suite covers reasoning (ARC Multilingual, CITE), math (MGSM CoT, CITE), logical reasoning (MLogiQA, CITE), instruction following (MIFEval, CITE), regional knowledge (Global MMLU, CITE; INCLUDE, CITE), and truthfulness (TruthfulQA, CITE).
%- arc multilingual (reasoning, MCQA)
%- global\_mmlu\_gen\_0shot (regional knowledge, MCQA)
%- mgsm\_en\_cot (math, generative)
%- include\_base\_44\_gen\_0shot (regional knowledge, MCQA)
%- mlogiqa\_gen
%- multi-if (instruction following, gen)
%- truthfulqa\_multilingual\_mc2 (MCQA)
%Given the scarcity of post-training multilingual benchmarks, we will also consider English-centric benchmarks like ... 


\section{Analysis}
\label{sec:results}

We organize the analysis around four questions, namely which benchmarks provide signal on small multilingual models (\ref{subsec:exp-da}), which SNR definition best predicts decision accuracy (\ref{subsec:exp-snr-da}), whether the resulting ranking generalizes across seeds and model suites (\ref{subsec:exp-generalization}), and whether benchmark subsets and design choices can improve reliability (\ref{subsec:exp-benchmark-design}).

\subsection{Calculation of Decision Accuracy}
\label{subsec:exp-da}

\textbf{Checkpoint selection.} For each custom and open-source model, we evaluate 10 evenly spaced checkpoints across training and 5 additional checkpoints in the final 10\% of training FLOPS. This schedule captures both the long-range trajectory and the late-training variation.

\textbf{The above-random gate.} Before computing DA or SNR, we filter (benchmark, size) cells whose mean score does not exceed the random baseline by at least 5 points. Figure~\ref{fig:gate} illustrates the two regimes. Of 118 benchmarks\footnote{Since we analyze results by language, we consider each language subset as a different benchmark (i.e. \texttt{hellaswag\_es} and \texttt{hellaswag\_hi} are two different benchmarks.}, only 44 clear chance at one or more custom model sizes, and 74 are at chance everywhere. The failures are almost entirely an answer-count effect, since four-option translated knowledge tasks sit at chance across all custom sizes while two-option tasks clear the gate. Notably, the three benchmarks that separate the data mixtures most are exactly the ones the gate removes. Their apparent signal is variation around chance, not capability. This shows that raw score separation cannot guide benchmark selection on its own.

\begin{figure}[t]
\centering
\includegraphics[width=0.7\textwidth]{figures/fig1_gate.pdf}
\caption{The above-random gate. Accuracy against training compute for the three data mixtures (colored curves) with final checkpoints marked. Left, a benchmark that clears the gate (\texttt{multiblimp}) shows scores well above the random baseline (dashed line) and consistent mixture separation.
Right, a gated benchmark (\texttt{belebele}) fluctuates around its baseline, so any apparent separation between mixtures is noise. A (benchmark, size) cell is kept iff its mean score exceeds the baseline by 5 points.}
\label{fig:gate}
\end{figure}

\textbf{Extension with external models.} Many tasks are at chance for the custom models, which are trained on 200 languages with only 100B tokens and at most 1B parameters. We therefore extend the analysis with 68 external models from 270M to 70B parameters (Appendix Table~\ref{tab:external}).
External sizes pool into buckets so that each bucket holds at least two models, and buckets contribute to signal and decision accuracy while singleton sizes only appear on the accuracy curves. With capable models in the pool, additional benchmarks clear the gate and enter the analysis.

% TODO: Figure or table showing acc vs flops and or top DA benchmarks, per DA type, and per language

\textbf{Per-language reliability at small scale.} Figure~\ref{fig:reliability-map} shows the full benchmark-by-language map for the above-random benchmarks on the small pretrained models. MultiBLiMP is the most reliable benchmark in six of eleven languages, with decision accuracy between 0.79 and 0.87, which matches the expectation that grammatical acceptability emerges early and therefore discriminates small models. Swahili retains no above-random benchmark on the custom models, so no reliable small-scale decision can be made for it at this scale. This confirms that benchmark reliability is language dependent and motivates the per-language reliability map as the practical output of the framework.

\textbf{Per-language reliability at 32B.} 
Table~\ref{tab:reliability-map} reports,
for each language, the benchmark whose small-model rankings best transfer to large models on the external suite. For each benchmark and language we average decision accuracy over all small-to-large size pairs, with small models up to 1.7B and large models from 7B to 32B parameters. Human-translated narrative and commonsense tasks dominate. XStoryCloze is the most transferable benchmark in four languages and close behind in three more, with decision accuracy up to 0.99, while machine-translated knowledge tasks rarely lead. The most reliable benchmark differs by language, and for Swahili no benchmark exceeds 0.57, barely above a random decision. This confirms that benchmark reliability is language dependent and that a single universal benchmark choice leaves some languages effectively unmeasured.

\begin{figure}[t]
\centering
\includegraphics[width=0.7\textwidth]{figures/fig3_reliability_map.pdf}
\caption{Benchmark reliability map. Decision accuracy by checkpoint at the 1B target for every benchmark family and language. Grey cells fail the above-random gate. Practitioners can read off which benchmarks support small-scale decisions for their target languages.}
\label{fig:reliability-map}
\end{figure}

\begin{table}[t]
\centering
\small
\caption{Per-language reliability map on the external suite (68 open-source models, 270M to 70B). For each language, the benchmark with the highest decision accuracy averaged over the nine small-to-large size pairs (small $\leq$ 1.7B, large 7B to 32B). DA is the fraction of model pairs whose small-scale ranking holds at large scale.}
\label{tab:reliability-map}
\begin{tabular}{llclc}
\toprule
Lang & Top benchmark & DA & Runner-up & DA \\
\midrule
en & \texttt{xwinograd}         & 0.81 & \texttt{xstorycloze} & 0.76 \\
es & \texttt{hellaswag}         & 0.97 & \texttt{xstorycloze} & 0.76 \\
ru & \texttt{xstorycloze}       & 0.99 & \texttt{hellaswag}   & 0.95 \\
hi & \texttt{xstorycloze}       & 0.78 & \texttt{multiblimp}  & 0.74 \\
zh & \texttt{xstorycloze}       & 0.97 & \texttt{xcopa}       & 0.76 \\
ja & \texttt{paws}              & 0.90 & \texttt{xwinograd}   & 0.83 \\
ar & \texttt{xstorycloze}       & 0.98 & \texttt{xnli}        & 0.81 \\
vi & \texttt{xcopa}             & 0.93 & \texttt{hellaswag}   & 0.73 \\
tr & \texttt{global\_mmlu\_full}& 0.92 & \texttt{xnli}        & 0.66 \\
th & \texttt{belebele}          & 0.90 & \texttt{xnli}        & 0.34 \\
sw & \texttt{belebele}          & 0.57 & \texttt{xcopa}       & 0.56 \\
eu & \texttt{hellaswag}         & 0.90 & \texttt{truthfulqa}  & 0.71 \\
\bottomrule
\end{tabular}
\end{table}

\subsection{Definition of SNR}
\label{subsec:exp-snr-da}

\textbf{SNR computation.} We compute the signal-to-noise ratio under each of the signal and noise definitions introduced in Section \ref{sec:methodology}, for every benchmark and every language. This lets us identify which definition correlates most strongly with decision accuracy for each language.

\textbf{Which SNR definition predicts decision accuracy.}
Table reports the mean Pearson correlation between log SNR and decision accuracy across languages, on the external suite and on the custom pool. The two populations differ in what signal measures, spread across controlled data mixtures for the custom models and spread across heterogeneous open-source models for the external suite, yet the result replicates. Under DA by checkpoint, the dispersion and relative-spread families reach correlations of 0.43 to 0.45 on the external suite and 0.43 to 0.51 on the custom pool, while depth and projection variants never help. Under DA by size the correlation is weak in both pools, for structural reasons. External size pairs compare different model families that differ in data and tokenizer, not only in size, and the custom 1B target is not fully converged. We therefore base benchmark selection on DA by checkpoint, and we recommend a metric family rather than a single variant.

\begin{figure}[t]
\centering
\includegraphics[width=0.5\textwidth]{figures/fig2_snr_vs_da.pdf}
\caption{SNR against decision accuracy by checkpoint at the 1B target, one
point per above-random benchmark, colored by language. Higher SNR predicts
higher decision accuracy.
}
\label{fig:snr-da}
\end{figure}

%\textbf{Language-dependent metric.} The same benchmark may be reliable for some languages and noisy for others. We therefore report all SNR and decision accuracy values per language and aggregate into a benchmark reliability map that practitioners can consult for the languages they target.

\subsection{Framework Generalization}
\label{subsec:exp-generalization}

\textbf{Generalization across seeds.} We test whether the variant ranking survives a change of training seed in our custom pretrained moddels by fitting on two seeds and holding out the third. The variant ranking transfers under DA by checkpoint, with a holdout Spearman correlation of 0.81, and does not transfer under DA by size, with a Spearman correlation of -0.07. The DA-size correlations are small and clustered within roughly 0.1 of each other, so their ranking is noise dominated. The exact best variant per language never transfers either. We therefore recommend a metric family rather than a single variant, and we treat DA-by-size variant rankings as unreliable at this scale.

\textbf{Agreement with the English SNR framework.} We compare our SNR values against those of \citet{heineman_signal_2025} on the DataDecide corpus. Seven English benchmarks are shared between the two model suites, of which four survive the above-random gate. Over these four tasks, SNR values computed on our custom models and on DataDecide agree closely, with a Pearson correlation of 0.98 on log SNR and a Spearman rank correlation of 1.00 under the \texttt{dispersion\_shifted} variant. With so few shared tasks, this is indicative rather than conclusive, but it suggests the framework measures a property of the benchmark rather than of the model suite. % A caveat applies to MMLU, since our suite runs the Cohere translation of the English split while DataDecide runs the original, so this row is not strictly like for like. We plan to rerun the original MMLU task on our checkpoints to remove the alias.

TODO: We will compare our benchmark selection and its compute cost against FineTasks \citep{kydlicek_finetasks_nodate} and SMART filtering \citep{gupta_improving_2025}.


\subsection{Improvement of Benchmark Quality}
\label{subsec:exp-benchmark-design}

%\textbf{Subsampling methods.} We explore subsampling rules that increase SNR on existing benchmarks, including selecting harder questions and removing semantically similar items. The output is a set of higher-SNR subsets that practitioners can use as replacements for full benchmark runs when running ablations. 

\textbf{Benchmarks subtasks are redundant.} Per benchmark, we rank
subtasks by standalone SNR and sweep cumulative subsets
(Figure~\ref{fig:subsets}). A subset of one to two MMLU subjects matches or beats the full set of roughly 48 subjects across model sizes, raising SNR from 2.12 to 3.65 at 175M, and the same subjects recur across sizes and languages. We do not read this as a recommendation to evaluate on a handful of subjects. We read it as evidence that a large share of benchmark content adds variance without adding discriminative signal, so the size of a benchmark is a poor proxy for its reliability. This invites a closer look at which portions of standard suites are informative and which are redundant.

\begin{figure}[t]
\centering
\includegraphics[width=0.6\textwidth]{figures/fig4_subset_sweep.pdf}
\caption{Cumulative subset sweep for Global MMLU subjects. Combined SNR as
subjects are added in order of standalone SNR, one line per model size, with
the full-set SNR as horizontal dashed lines. Small subsets match or beat the
full set.
% TODO: confirm best_n per size from regenerated figure. 
Per-item selection yields larger gains
that do not transfer across sizes.}
\label{fig:subsets}
\end{figure}

\textbf{Item-level subsets.} The same procedure at the level of individual items finds subsets with even larger SNR gains, showing that informative items exist and can be identified within a single model size. The selected items, however, differ across model sizes, with a Jaccard overlap of 0.03 between the best subsets at different scales. Item-level selection is therefore promising as a per-scale tool, and making the selected subsets transfer across scales is the main open problem, which we continue to investigate.

TODO: We will pursue item-level subset selection that transfers across scales, combining our SNR ranking with item response theory approaches \citep{zhou_lost_2026}.


\textbf{Benchmark creation decisions.} We document the data source and curation method of each benchmark in our suite, and compute the correlation between these characteristics and SNR. We classify the data source as web data, news, literature, crowdsourced, expert-generated, or synthetic, and the curation method by whether the data was machine or manually translated and by the review process applied.
We test whether benchmark design features predict SNR on the external suite, where eleven benchmark families clear the above-random gate, including four-option knowledge tasks that capable models solve. No single feature is significant. Family-level Kruskal-Wallis tests on curation method ($p = 0.49$), task format ($p = 1.00$), answer-option count ($p = 0.83$), and source origin ($p = 0.68$) are all far from significance. How a benchmark was built does not predict its reliability once models clear chance on it. The pattern that does hold sits in the answer space, since every benchmark that fails the gate on small models is a four-option translated knowledge task, and benchmarks are sharper when the model compares fewer and longer scored completions. This suggests that for small-scale evaluation, reducing the answer space matters more than curation quality.

%\input{sections/04_analysis_b} % RQ4 & 5
\section{Conclusions}
\label{sec:conclusions}

We study when multilingual evaluation supports decisions during pretraining. Predictable learning curves, evidence above chance, and agreement with a larger model answer different questions. HellaSwag, LAMBADA, and per-language BPB track scale consistently, while reformulating letter-based tasks as answer continuations reveals signal that their original protocols suppress. Across the 12-checkpoint evaluation grid, BPB rankings transfer more consistently to the 1.7B reference than the aggregate benchmark rankings, but early agreement can reverse during training. Cross-task transfer is limited, and the most informative SNR definition varies by language and by whether the decision concerns model size or training progress. These findings favor task- and language-specific checks over a single benchmark score or universal proxy threshold.

Our next steps address two questions left open by this analysis. First, we will test how predictivity depends on the design intervention, comparing language selection, sampling temperature, second-language choice, and architecture while estimating uncertainty in the reference ranking from repeated runs. Second, we will test whether scaling behavior and decision reliability transfer to held-out languages, with explicit separation between languages absent from pretraining and languages withheld from fitting the predictor. These studies will also provide held-out decisions for calibrating surrogate-based selection and stopping rules, extending the present descriptive evidence toward actionable guidance for multilingual model development.


\bibliography{zotero,custom}
\bibliographystyle{iclr2027_conference}

\appendix
\section{Terminology}

\begin{itemize}
    \item Benchmark \(b\): HellaSwag.
    \item Language task \(b_\ell\): HellaSwag-es, HellaSwag-en, HellaSwag-fr.
    \item Sample/example: one question or instance inside HellaSwag-es.
    \item Evaluation suite: the complete collection of benchmark families and tasks.
\end{itemize}

\begin{table}[h]
    \centering
    \caption{Notation symbols}
    \label{tab:notation-symbols}
    \begin{tabular}{lll}
        \hline
        \textbf{Symbol} & \textbf{Meaning} & \textbf{Example} \\
        \hline
        \(N\) & Non-embedding parameter count & 350M \\
        \hline
        \(A\) & Architecture & deep \\
        \hline
        \(K\) & Number of training languages & 15 \\
        \hline
        \(M\) & Language-selection scheme or mixture & scheme A \\
        \hline
        \(T\) & Sampling temperature & 1 or 3 \\
        \hline
        \(S\) & Random seed & 1904 \\
        \hline
        \(f\) & Fraction of training completed & 0.20 \\
        \hline
        \(b\) & Benchmark family & HellaSwag \\
        \hline
        \(\ell\) & Evaluation language & Spanish \\
        \hline
        \(b_\ell\) & Language task & HellaSwag-es \\
        \hline
        \(Y_{b_\ell}\) & Raw task score & accuracy or BPB \\
        \hline
    \end{tabular}
\end{table}


% Requires: \usepackage{booktabs}
\begin{table}[h]
    \caption{Terminology}
    \label{tab:terminology}
    \centering
    \begin{tabular}{lll}
        \toprule
        Term & Use it for & Do not use it for \\
        \midrule
        \textbf{Model configuration} & $\theta$: \((N,A,K,M,T)\) & A seed replicate \\
        \textbf{Training run} & One configuration trained with seed \(S\) & One task evaluation \\
        \textbf{Checkpoint} & Saved model state at fraction \(f\) & The fraction itself \\
        \textbf{Training fraction} & \(f\in\{0.1,\ldots,1.0,0.85,0.95\}\) & A model artifact \\
        \textbf{Model variant} & One level of an intervention being ranked & A proxy/reference checkpoint \\
        \textbf{Proxy point} & Size and training fraction used for the inexpensive ranking & The model variant \\
        \textbf{Reference point} & Size and fraction defining the target ranking & Necessarily 1.7B \\
        \textbf{Benchmark} & HellaSwag & HellaSwag-es \\
        \textbf{Language task} & HellaSwag-es & The complete HellaSwag family \\
        \textbf{Sample/example} & One question & A task \\
        \textbf{Score observation} & One checkpoint evaluated on one task & ``Evaluation run'' \\
        \textbf{Ranking agreement} & General concept & A numerical metric \\
        \textbf{Decision accuracy} & Pairwise ranking-agreement metric & General reliability \\
        \textbf{Scaling predictability} & \(R^2\) of a scaling fit & Decision accuracy \\
        \textbf{Scaling monotonicity} & Spearman correlation along a scaling curve & Rank transfer \\
        \textbf{Reliability} & Broad interpretive concept & A metric unless explicitly defined \\
        \bottomrule
    \end{tabular}
\end{table}

\begin{table}[h]
    \caption{Notation $\rho$}
    \label{tab:notation-rho}
    \centering
    \begin{tabular}{|l|l|}
        \hline
        Quantity & Recommended notation \\ \hline
        Proxy/reference rank agreement & \(\rho_{\mathrm{RA}}\) \\ \hline
        Scaling monotonicity & \(\rho_{\mathrm{scale}}\) \\ \hline
        Association between surrogate and DA & \(\rho_{\mathrm{assoc}}\) \\ \hline
    \end{tabular}
\end{table}

% =============================================================================
% Methodology and Experimental Setup — the predictivity sweep
%
% Sources (repo, as of 2026-09-25):
%   plan/small-to-large-predictivity-training-plan.md   design, budgets, data, validation
%   plan/methodology-review-2026-08-21.md               LR law, shallow ladder, checkpoint grid
%   plan/benchmark_selection.md                         benchmark families + per-cell task counts
%   plan/90M-rung-anomaly.md                            the 90M divergence (appendix draft)
%   plan/1b-models.md                                   the 20-checkpoint 1B cells
%   src/pretrain/launch_trainings.py                    DATA_SCHEMES, SEED_TRIPLES, n_checkpoints
%   src/pretrain/hyperparams/hyperparams_{deep,shallow}.json
%   src/pretrain/megatron_args.sh                       every training argument
%   src/pretrain/data/{create_data_mixture,build_data_mixtures}.py, language_sets_scheme*.json
%   src/evals/scripts/score_bpb.py                      BPB
%   src/evals/README.md, src/pretrain/auto_evals_cscs.py harness settings
%   configs/tasks.json                                  the `auto` group and family metadata
%
% Included in the appendix through app_languages.tex. Citation keys follow the
% .bib's Zotero style; keys present in the other sections are reused verbatim,
% the following are GUESSED and must be checked against the .bib:
%   ainslie_gqa_2023 hoffmann_training_2022 bhagia_establishing_2024
%   pagliardini_ademamix_2024 hagele_scaling_2024 gao_framework_2024
%   kwon_efficient_2023 magnusson_datadecide_2025
% Anything still marked \citeneeded{...} has no key yet.
% =============================================================================
\providecommand{\citeneeded}[1]{\textcolor{red}{[CITE: #1]}}
\providecommand{\checknum}[1]{\textcolor{orange}{#1}}


\section{Model Ladder Details}
\label{app:methodology}

We give the training grid and checkpoint schedule (\S\ref{app:training-grid}), model configurations (\S\ref{app:model-configurations}), data construction (\S\ref{app:data-construction}), and evaluation protocol (\S\ref{app:evaluation-protocol}).


\subsection{Training Grid and Checkpoints}
\label{app:training-grid}

\begin{table}[t]
\centering
\small
\setlength{\tabcolsep}{4pt}
\caption{The data schemes of the sweep. Every scheme trains the six-rung ladder from 90M to 1.7B at the settings it defines; the 3B rung is trained only under schemes A and B at $L \in \{8, 15\}$, deep, and is counted in the A and B rows. Seeds: ``grid'' means one seed per cell, or three seeds for the deep cells marked in Table~\ref{tab:grid}; every other scheme trains a single seed. Runs count both architectures where the scheme trains both. A further scheme, BT3 (the scheme-B $L=30$ list at $T=3$), is registered and built but not trained.}
\label{tab:schemes}
\begin{tabular}{llccllr}
\toprule
Scheme & Language lists & $T$ & $L$ & Sizes & Architectures & Runs \\
\midrule
A   & resource-ranked (baseline) & 1 & 1, 2, 8, 15, 30, 50 & all (+3B) & deep, shallow & 92 \\
AT3 & scheme A lists, flattened  & 3 & 15, 30, 50          & all & deep, shallow ($L=50$); deep ($L=15, 30$) & 24 \\
B   & diversity-first            & 1 & 8, 15, 30           & all (+3B) & deep, shallow & 40 \\
ZH  & Chinese as the second language & 1 & 2              & all & deep     & 6 \\
ES  & Spanish as the second language & 1 & 2              & all & deep     & 6 \\
\midrule
Total & & & & & & 168 \\
\bottomrule
\end{tabular}
\end{table}

\textbf{Seeds.} Every cell trains with seed 1904. To measure how much a ranking can move for reasons unrelated to the intervention, a subset of cells trains three seeds, each with a different initialization and data order (Table~\ref{tab:grid}): seeds 64, 313, and 1904 at 175M and 600M for $L \in \{1, 2, 50\}$, and seeds 28, 1797, and 1904 at 1B for $L \in \{1, 2, 30\}$, all in the deep architecture. The seed triples cover 175M, 600M, and 1B; they do not directly estimate seed variability at 1.7B. The extra schemes (AT3, ZH, ES) are intervention levels rather than part of the noise estimate and train a single seed.

\begin{table}[t]
\centering
\small
\caption{The scheme-A deep grid: seeds per $(N, L)$ cell. The shallow architecture trains one seed per cell. The 3B rung (not shown) trains $L \in \{8, 15\}$ with one seed. In total the scheme-A deep grid has 56 runs, the 90M and 3B rungs included.}
\label{tab:grid}
\begin{tabular}{lcccccc}
\toprule
$L$ & 90M & 175M & 350M & 600M & 1B & 1.7B \\
\midrule
1   & 1 & 3 & 1 & 3 & 3 & 1 \\
2   & 1 & 3 & 1 & 3 & 3 & 1 \\
8   & 1 & 1 & 1 & 1 & 1 & 1 \\
15  & 1 & 1 & 1 & 1 & 1 & 1 \\
30  & 1 & 1 & 1 & 1 & 3 & 1 \\
50  & 1 & 3 & 1 & 3 & 1 & 1 \\
\bottomrule
\end{tabular}
\end{table}

\textbf{Checkpoints.} Each run saves evenly spaced checkpoints: 20 per run, 40 at the 1B rung, and 60 at the 1.7B rung, so that the two reference rungs are sampled more densely. Because 40 and 60 are multiples of 20, the denser grids contain the shared fractions $k/20$, for $k=1,\ldots,20$. The $20N$ Chinchilla token budget falls at one fifth of training. We evaluate twelve checkpoints per run: the ten tenths of training, which every size shares and which let us compare sizes at the same multiple of Chinchilla (20\% of a run is 1C and its end 5C), plus the 85\% and 95\% saves. The analysis reads the ten tenths. The two extra points exist so that checkpoint-to-checkpoint noise can be estimated over five evenly spaced checkpoints in the last 20\% of training (80, 85, 90, 95 and 100\%), the decay phase of the schedule.\footnote{Fifteen 1B cells were trained before the 40-checkpoint rule and save 20 checkpoints (every 2,287 steps to step 45,740 rather than every 1,143 steps to 45,720). Their saves land on the same $k/20$ fractions, so they are evaluated at the same points as every other run; the token count at matched checkpoints differs by 0.044\%.}

\textbf{Compute axis.} Every checkpoint has a known $(N, D)$, so results can also be read against training compute. We estimate $\mathrm{FLOPs} = 6\,(N + d_{\mathrm{model}} V)\,D$, where $V$ is the vocabulary size: we omit embedding-lookup cost but include the output projection, and with a 131K vocabulary the projection nearly doubles the per-token cost of the smallest rungs.

\subsection{Model Configurations and Optimization}
\label{app:model-configurations}

\textbf{Parameter-count convention.} All sizes are non-embedding parameters, that is, the transformer blocks alone, excluding the tied token embedding and output projection. This follows the ladder convention of \citep{bhagia_establishing_2024} adopted by \citet{heineman_signal_2025}. The convention matters here because the 131,072-token vocabulary makes the tied embedding matrix $d_{\mathrm{model}} \times V$ parameters on its own, 101M at the 90M rung and 302M at the 1.7B rung; including it would substantially change the relative spacing of the smallest models on the size axis.

\textbf{Architecture.} All models are decoder-only transformers following the Apertus pretraining recipe \citep{apertus_apertus_2025}: grouped-query attention \citep{ainslie_gqa_2023} with four query heads per key-value head and a head dimension of 64, rotary position embeddings \citep{su_roformer_2023} with base 500,000 and Llama-3-style scaling factor 8, RMSNorm, QK-layernorm, the xIELU activation \citep{huang_xielu_2025}, no bias terms, no dropout, a sequence length of 4,096, and tied input and output embeddings over the 131,072-token Apertus tokenizer (\texttt{swiss-ai/Apertus-70B-2509}). The deep ladder scales by jointly increasing depth, width, and the number of attention heads, keeping a width-to-depth ratio $d_{\mathrm{model}}/n_{\mathrm{layers}}$ of 64 from 175M to 1B (51 at 90M and 77 at 1.7B, the endpoint compromises needed to hit the target sizes) and a feed-forward multiplier of 4. Table~\ref{tab:ladder} gives the configurations.

\textbf{Shallow ladder.} The depth intervention trains a second ladder at the same six non-embedding sizes with roughly twice the width-to-depth ratio (110 to 146 against 51 to 77). To constrain the architecture comparison, the shallow search pins the feed-forward multiplier and the query-to-key-value head ratio to the deep values, restricts $d_{\mathrm{model}}$ to multiples of 256, and among candidates within 4\% of the target size picks the ratio closest to 128. The shallow sizes land within 0.4\% to 5.2\% of the deep ones, and each shallow run trains its own $D = 100N$.

\begin{table}[t]
\centering
\small
\setlength{\tabcolsep}{3.5pt}
\caption{The model ladder. $N$ counts non-embedding parameters; the tied embedding adds $d_{\mathrm{model}} \times 131{,}072$. $D = 100N$ is the token budget; steps at 504 sequences of 4,096 tokens (2.06M tokens per step); LR is the peak learning rate from the compute-based law at the run's own budget; ckpts is the number of saved checkpoints. GQA (grouped-query attention) gives query/key-value head counts. Shallow rows approximately match the deep models in $N$; their token budgets follow their own parameter counts. FFN is the feed-forward hidden dimension.}
\label{tab:ladder}
\begin{tabular}{llrrrrrrrrr}
\toprule
Arch & Size & Layers & $d_{\mathrm{model}}$ & FFN & GQA & $N$ (M) & $D$ (B) & Steps & LR ($10^{-3}$) & Ckpts \\
\midrule
deep & 90M  & 15 &  768 & 3,072 & 12/3  &   92.9 &   9.3 &  4,500 & 1.43 & 20 \\
deep & 175M & 16 & 1,024 & 4,096 & 16/4  &  176.2 &  17.6 &  8,540 & 1.22 & 20 \\
deep & 350M & 20 & 1,280 & 5,120 & 20/5  &  344.1 &  34.4 & 16,660 & 1.03 & 20 \\
deep & 600M & 24 & 1,536 & 6,144 & 24/6  &  594.5 &  59.5 & 28,800 & 0.90 & 20 \\
deep & 1B   & 28 & 1,792 & 7,168 & 28/7  &  944.1 &  94.4 & 45,720 & 0.80 & 40 \\
deep & 1.7B & 30 & 2,304 & 9,216 & 36/9  & 1,672.2 & 167.2 & 81,000 & 0.69 & 60 \\
\midrule
shallow & 90M  &  8 & 1,024 &  4,096 & 16/4  &   88.1 &   8.8 &  4,260 & 1.45 & 20 \\
shallow & 175M & 10 & 1,280 &  5,120 & 20/5  &  172.0 &  17.2 &  8,340 & 1.22 & 20 \\
shallow & 350M & 14 & 1,536 &  6,144 & 24/6  &  346.8 &  34.7 & 16,800 & 1.03 & 20 \\
shallow & 600M & 14 & 2,048 &  8,192 & 32/8  &  616.6 &  61.7 & 29,860 & 0.89 & 20 \\
shallow & 1B   & 17 & 2,304 &  9,216 & 36/9  &  947.6 &  94.8 & 45,920 & 0.80 & 40 \\
shallow & 1.7B & 20 & 2,816 & 11,264 & 44/11 & 1,665.3 & 166.5 & 80,640 & 0.69 & 60 \\
\bottomrule
\end{tabular}
\end{table}

\textbf{Optimization.} Every run uses the same optimizer and schedule, with only the per-size values changing. We train with AdEMAMix \citep{pagliardini_ademamix_2024} ($\beta_1 = 0.9$, $\beta_2 = 0.999$, $\beta_3 = 0.9999$, $\alpha = 8$), weight decay 0.1, gradient clipping at 0.1, and a global batch of 504 sequences of 4,096 tokens, about 2.06M tokens per step, at every size. The learning rate follows a warmup-stable-decay schedule \citep{hagele_scaling_2024}: linear warmup over about 4\% of the steps, a constant phase, and a $1-\sqrt{\cdot}$ decay to zero over the final 20\%. The peak learning rate comes from a compute-based law evaluated at each run's own budget, $\eta = 0.3118\, C^{-1/8}$ with $C = 6 N D = 600 N^2$, so it decreases smoothly from $1.43 \times 10^{-3}$ at 90M to $6.9 \times 10^{-4}$ at 1.7B. We scale optimizer warmup with run length and initialization with width: the AdEMAMix $\alpha$ and $\beta_3$ warmup spans the full run, and the initialization standard deviation is width-scaled as $0.008944 \sqrt{1792 / d_{\mathrm{model}}}$. Training uses BF16 mixed precision with FP32 gradient accumulation in the Swiss AI fork of Megatron-LM \citep{shoeybi_megatron-lm_2019}, with data parallelism only (no tensor or pipeline parallelism).

\textbf{The 90M rung.} Nine of the ten completed 90M runs diverge: each reaches its best training loss between 15\% and 19\% of the way through training and degrades for the rest of the run, ending 1.2 to 1.9 nats above its own best. No larger rung shows this. Held-out bits-per-byte rises with training on every language we checked, so the degradation is real rather than an artefact of the training objective. Retraining the rung at three learning rates spanning a factor of five and at one half and one sixth of the batch size reproduces the failure at the same fraction of the optimizer steps, which points to the interaction between the fixed $\beta_3$ timescale of AdEMAMix, about 10,000 steps, and a run that is only 4,500 steps long. We keep the rung in the grid and report its results, but exclude it from the scaling fits. Appendix~\ref{app:90m-rung} documents the investigation.


\subsection{Data Construction}
\label{app:data-construction}

\textbf{Language lists.} The language lists are generated from the FineWeb-2 per-language byte counts and the benchmark availability in the evaluation harness. \emph{Scheme A} ranks FineWeb-2 subsets by training-split UTF-8 bytes, excluding the unidentified-language subsets, and takes the top $L-1$; the lists are nested, so the $L=2$ list (Russian) is a subset of the $L=8$ list (Russian, Chinese, German, Japanese, Spanish, French, Italian), which is a subset of the $L=15$, $L=30$, and $L=50$ lists. A 100-language list was generated by the same rule (with the eight subsets among the top 99 that have no benchmark in the harness replaced by the next subsets by bytes with at least two benchmark families) but was not trained, as explained below. \emph{Scheme B} builds diversity-first lists at $L \in \{8, 15, 30\}$ within the top 49: it takes the largest subset of each script not yet covered, then of each language family not yet covered, then the remaining subsets by bytes. At $L=8$ the scheme-B list spans seven scripts against four for scheme A, and at $L=30$ it spans 14 scripts and 12 families against 10 and 10. \emph{Schemes ZH and ES} replace Russian with Chinese or Spanish at $L=2$. The full lists are given in Appendix~\ref{app:languages}.

\textbf{Mixture composition.} Every multilingual setting blends English and FineWeb-2 50/50 at training time with the Megatron data loader's blend weights, from one English dataset and one FineWeb-2 dataset per (scheme, setting). Within the FineWeb-2 half, language $i$ receives a share $q_i \propto p_i^{1/T}$, where $p_i$ is its estimated token proportion in the corpus. Scheme A allocates at $T=1$, proportionally to the corpus. The FineWeb-2 byte distribution is extreme: on the planned 100-language list at $T=1$ the largest language would take 28\% of the multilingual half and the smallest 0.0017\%, and at the 90M rung 66 of the 99 languages would receive under 10M tokens. The flattened scheme AT3 allocates at $T=3$, which would lift the median language of that build from 90M to 373M tokens. The $L=100$ setting was dropped from the grid; the temperature intervention is instead trained on the scheme-A lists at $L=50$ (both architectures) and, where flattening starves no language, at $L=30$ and $L=15$ (deep only), whose smallest language keeps 343M and 1.2B tokens at $T=1$ under the 1.7B budget.

\textbf{Build sizing and repetition.} Each FineWeb-2 dataset is built to half of the largest budget that trains on it plus 10\% headroom, 92B tokens for the 1.7B rung; the ZH build (52B) was sized when that scheme stopped at 1B and is read unchanged by its 1.7B cell. The English dataset is built once to 184B tokens. The builder never repeats a document: a language that runs out of data yields a smaller build. Smaller models draw a shuffled fraction of the same build, so within a setting all sizes see the same language proportions. Two departures from a single pass are recorded and carried into the analysis. First, no single FineWeb-2 subset can feed the 1.7B rung at $L=2$: the Russian build realizes 72.8B tokens against the 83.6B a 1.7B draws from the multilingual half, so the 1.7B at $L=2$ sees Russian about 1.15 times. The Chinese build (52.0B tokens) and the Spanish build (23.9B, all the Spanish the source holds) are smaller still: at 1.7B, Chinese is seen 1.61 times and Spanish 3.51 times (0.65, 0.91 and 1.98 times for Russian, Chinese and Spanish at 1B). Repetition stays below 1.7 epochs across the Russian and Chinese ladders but rises from 0.7 at 350M to 3.5 at 1.7B for Spanish, so any comparison of the $L=2$ schemes states these epoch counts. Second, the $L=15$ and $L=50$ builds of scheme A and the $L=15$ build of scheme B were originally built to 52B and extended to 92B when the 1.7B rung was added at those settings; the extension continues each language's token stream from the same starting point, so the six 1.7B cells there read a superset of what their proxies read, extended with later CommonCrawl dumps.\footnote{Verified byte for byte: every language section of each 52B build is a prefix of the corresponding 92B build.}

\textbf{Validation set.} One fixed validation set is shared by every model and every scheme. For each of the 99 FineWeb-2 languages of the $L=100$ list plus English, it takes the first 5M tokens of that language's first parquet file, capped at 30\% of the file's rows so that single-file languages keep training data. The build records each language's token count, its UTF-8 byte count, and the number of leading rows assigned to validation; every training build reads this manifest and skips exactly those rows, so training and validation are disjoint by construction. Since there is one validation build for the whole sweep, bits-per-byte is comparable across sizes, language settings, and schemes.

\subsection{Evaluation Protocol}
\label{app:evaluation-protocol}

\textbf{Bits-per-byte implementation.} We measure the negative log-likelihood of held-out text in bits, normalized by its UTF-8 byte count.
For language $\ell$, let $x_1,\ldots,x_{m_\ell}$ be the evaluated token sequence, $c_t$ the context preceding token $x_t$ within the scoring block, and $b_\ell$ the number of bytes in the evaluated text. We compute
\begin{equation}
\mathrm{BPB}_\ell = \frac{-\sum_{t=2}^{m_\ell}\log_2 p(x_t \mid c_t)}{b_\ell}.
\label{eq:bpb}
\end{equation}

We concatenate validation documents without end-of-document separators and score them in blocks with 4,096 input tokens and one-token overlap between successive blocks. Each token except the first of the stream is scored once; context resets at block boundaries. We select a deterministic prefix ending at the document boundary that reaches 1M tokens, or the full stream if it is shorter. We measure the denominator by decoding the entire selected prefix and counting its UTF-8 bytes, including the initial conditioning token. Thus, the implementation omits the first token's likelihood but retains its bytes. We also compute token-level perplexity from the same forward passes.

% The original draft reports a standard error of about 0.003 bits for this
% prefix. Add the estimation procedure before presenting it as uncertainty.

\begin{table}[t]
\centering
\small
\setlength{\tabcolsep}{4pt}
\caption{Benchmark families evaluated during pretraining. ``Langs'' is the number of languages the harness ships for the family; a model is evaluated only on the languages in its training mixture. MC = multiple choice; QA = question answering; UD = Universal Dependencies. Construction: HT = human or professional translation, MT = machine translation, native = written in the language, auto = automatically generated.}
\label{tab:benchmark-families}
\begin{tabular}{llllr}
\toprule
Family & Capability & Format & Construction & Langs \\
\midrule
Belebele          & reading comprehension     & 4-way MC, 900 items/lang      & HT (FLORES)           & 122 \\
Global PIQA           & physical commonsense      & 2-way completion              & native                & 116 \\
MultiBLiMP            & grammatical acceptability & minimal pairs                 & auto (UD, UniMorph)   & 101 \\
INCLUDE             & regional knowledge        & MC exam questions             & native                &  44 \\
Global-MMLU  & world knowledge & MC (MMLU)                & HT + community        &  42 \\
Okapi HellaSwag           & activity commonsense      & 4-way completion              & MT                    &  31 \\
Okapi ARC             & science QA                & MC                            & MT                    &  31 \\
XNLI                 & natural language inference & 3-way                        & HT                    &  15 \\
XStoryCloze              & narrative commonsense     & 2-way ending                  & HT                    &  11 \\
XCOPA                  & causal commonsense        & 2-way, 500 items              & HT                    &  11 \\
PAWS-X                 & paraphrase identification & 2-way                         & HT                    &   7 \\
XWinograd     & coreference               & Winograd schemas              & HT                    &   6 \\
LAMBADA multilingual & last-word prediction & completion             & MT                    &   5 \\
TruthfulQA-Multi  & truthfulness           & MC                            & HT                    &   5 \\
\bottomrule
\end{tabular}
\end{table}

\textbf{Task selection.} We select the registered tasks whose language is in the training mixture; task counts can exceed family counts. For the resource-ranked lists, the number of tasks grows with $L$: 106 at $L=1$, 125 at $L=2$, 294 at $L=8$, 430 at $L=15$, 590 at $L=30$, and 776 at $L=50$. These counts cover the whole evaluated group, which besides the multilingual families of Table~\ref{tab:benchmark-families} contains the reformulated variants of the letter-format families (\S\ref{subsec:rq1}) and a set of English-only benchmarks \checknum{(the table predates their addition and must be regenerated)}. Every trained language is covered by at least one family. For the seed-1904 deep scheme-A ladder we additionally run every task in every language, independent of the training mixture, to read cross-lingual transfer on the benchmarks as well as on bits-per-byte. Appendix~\ref{app:benchmark-coverage} gives the per-language coverage.

\textbf{Harness settings.} All benchmarks run through the Swiss AI fork of the LM Evaluation Harness \citep{gao_framework_2024}, installed from a pinned shared checkout with the source commit recorded in each evaluation job, with the vLLM backend \citep{kwon_efficient_2023}, a maximum context of 4,096 tokens, a beginning-of-sequence token prepended to every prompt, no chat template, and each task's default prompt and number of in-context examples. We report accuracy for multiple-choice and completion tasks and exact match for the generative ones, and aggregate sub-tasks (for example the MMLU subjects) to their family by the harness's own grouping. Base models at these sizes are at chance on several families; the above-chance gate of Appendix~\ref{app:results-filtering} decides per task and size whether such a cell enters the analysis.

\section{Evaluations filtering}
\label{app:results-filtering}

For each benchmark--language subset $(B,L)$ and model size $N$, we test whether the observed accuracy is distinguishable from random guessing.
For a task with $K$ answer options, chance accuracy is defined as $1/K$.
For every evaluation run, we compute the one-sided 95\% Wilson lower confidence bound on its accuracy and mark the run as above chance when

\[
\operatorname{LCB}_{95}(\hat{p}) > \frac{1}{K}.
\]

We then aggregate these decisions across the available architectures, language mixtures, and seeds for the same $(B,L,N)$ cell.
A cell is retained if at least 50\% of its evaluation runs are above chance.
When models trained on the benchmark language are available, we use those runs for the filtering decision; otherwise, we use all evaluated runs.
Measurements without a meaningful chance baseline, such as bits per byte and training loss, are not filtered by this procedure.

The Wilson bound accounts for the number of evaluation examples, so the amount by which accuracy must exceed chance varies across tasks.
Smaller evaluation sets require a larger margin above chance, while larger sets can support smaller but more reliable departures from chance.
For example, the median required margins are 9.00 percentage points for nonparallel Global PIQA, 2.44 for Belebele, 0.74 for HellaSwag, and 0.60 for Global-MMLU.
The corresponding fractions of benchmark--language--size cells that pass the filter are 37.5\%, 6.8\%, 92.7\%, and 4.3\%, respectively.
This filtering step is used only to determine whether a benchmark provides measurable signal at a given scale; it does not assess whether the benchmark can reliably rank model variants.

\section{Languages, benchmarks, and models}

\begin{table*}[ht]
\centering
\small
\setlength{\tabcolsep}{5pt}
\caption{Language selection. We cover 8 language families plus one isolate, 7 writing systems, and high to low-resource.}
\label{tab:languages}
\begin{tabular}{lllll}
\toprule
Code & Language & Family & Script & Resource \\
\midrule
en & English & Indo-European, Germanic & Latin & High \\
es & Spanish & Indo-European, Romance & Latin & High \\
ru & Russian & Indo-European, Slavic & Cyrillic & High \\
hi & Hindi & Indo-European, Indo-Aryan & Devanagari & Mid \\
zh & Mandarin Chinese & Sino-Tibetan & Han & High \\
ja & Japanese & Japonic & Kanji and Kana & High \\
ar & Modern Std. Arabic & Afro-Asiatic, Semitic & Arabic & High \\
vi & Vietnamese & Austroasiatic, Vietic & Latin & Mid \\
tr & Turkish & Turkic & Latin & Mid \\
th & Thai & Tai-Kadai & Thai & Mid \\
sw & Swahili & Niger-Congo, Bantu & Latin & Low \\
eu & Basque & Isolate & Latin & Low \\
\bottomrule
\end{tabular}
\end{table*}

\begin{table*}[ht]
\centering
\small
\setlength{\tabcolsep}{4pt}
\caption{Benchmark suite.}
\label{tab:benchmarks}
\begin{tabular}{llcccccccccccc}
\toprule
Benchmark & Capability & en & es & ru & hi & zh & ja & ar & vi & tr & th & sw & eu \\
\midrule
\texttt{arc} & scientific knowledge & \checkmark & \checkmark & \checkmark & \checkmark & \checkmark &  & \checkmark & \checkmark &  &  &  & \checkmark \\
\texttt{belebele} & reading comprehension & \checkmark & \checkmark & \checkmark & \checkmark & \checkmark & \checkmark & \checkmark & \checkmark & \checkmark & \checkmark & \checkmark & \checkmark \\
\texttt{global\_mmlu\_full} & world knowledge & \checkmark & \checkmark & \checkmark & \checkmark & \checkmark & \checkmark & \checkmark & \checkmark & \checkmark &  & \checkmark &  \\
\texttt{global\_piqa} & physical commonsense & \checkmark & \checkmark & \checkmark & \checkmark & \checkmark & \checkmark & \checkmark & \checkmark & \checkmark & \checkmark & \checkmark &  \\
\texttt{hellaswag} & activity commonsense & \checkmark & \checkmark & \checkmark & \checkmark &  &  & \checkmark & \checkmark &  &  &  & \checkmark \\
\texttt{include\_base\_44} & regional knowledge &  & \checkmark & \checkmark & \checkmark & \checkmark &  & \checkmark & \checkmark & \checkmark &  &  & \checkmark \\
\texttt{mlqa} & cross-lingual QA & \checkmark & \checkmark &  & \checkmark & \checkmark &  & \checkmark & \checkmark &  &  &  &  \\
\texttt{multiblimp} & grammatical acceptability & \checkmark & \checkmark & \checkmark & \checkmark &  &  & \checkmark &  & \checkmark &  &  & \checkmark \\
\texttt{paws} & paraphrase detection & \checkmark & \checkmark &  &  & \checkmark & \checkmark &  &  &  &  &  & \checkmark \\
\texttt{truthfulqa} & truthfulness &  & \checkmark & \checkmark & \checkmark & \checkmark &  & \checkmark & \checkmark &  &  &  & \checkmark \\
\texttt{xcopa} & causal commonsense &  &  &  &  & \checkmark &  &  & \checkmark & \checkmark & \checkmark & \checkmark & \checkmark \\
\texttt{xnli} & natural language inference & \checkmark & \checkmark & \checkmark & \checkmark & \checkmark &  & \checkmark & \checkmark & \checkmark & \checkmark & \checkmark & \checkmark \\
\texttt{xstorycloze} & narrative commonsense & \checkmark & \checkmark & \checkmark & \checkmark & \checkmark &  & \checkmark &  &  &  & \checkmark & \checkmark \\
\texttt{xwinograd} & coreference resolution & \checkmark &  & \checkmark &  & \checkmark & \checkmark &  &  &  &  &  &  \\
\texttt{mgsm\_direct} & mathematical reasoning & \checkmark & \checkmark & \checkmark &  & \checkmark & \checkmark &  &  &  & \checkmark & \checkmark & \checkmark \\
\texttt{mgsm\_native\_cot} & mathematical reasoning & \checkmark & \checkmark & \checkmark &  & \checkmark & \checkmark &  &  &  & \checkmark & \checkmark &  \\
\texttt{multi-if} & instruction following &  &  &  &  &  &  &  &  &  &  &  &  \\
\bottomrule
\end{tabular}
\end{table*}


\begin{table*}
\tiny
\centering
\caption{Coverage of the \texttt{external} model set: 68 open-source models.}
\begin{tabular}{llrrr}
\toprule
Model & Size & Ckpts & Tasks & Families \\
\midrule
\texttt{gemma-3-270m} & 270M & 1 & 1006 & 30 \\
\texttt{gemma-3-270m-it} & 270M & 1 & 1006 & 30 \\
\texttt{OLMo-2-0425-1B} & 1B & 1 & 1006 & 30 \\
\texttt{OLMo-2-0425-1B-DPO} & 1B & 1 & 1006 & 30 \\
\texttt{OLMo-2-0425-1B-Instruct} & 1B & 1 & 1006 & 30 \\
\texttt{OLMo-2-0425-1B-SFT} & 1B & 1 & 1006 & 30 \\
\texttt{apertus-1b-from8b-TOP256-long} & 1B & 10 & 1006 & 30 \\
\texttt{apertus3-1b-21-nodes} & 1B & 18 & 1006 & 30 \\
\texttt{gemma-3-1b-it} & 1B & 1 & 1021 & 30 \\
\texttt{gemma-3-1b-pt} & 1B & 1 & 1006 & 30 \\
\texttt{Qwen3-1.7B} & 1.7B & 1 & 1026 & 30 \\
\texttt{Qwen3-1.7B-Base} & 1.7B & 1 & 1006 & 30 \\
\texttt{SmolLM3-3B} & 3B & 1 & 1026 & 30 \\
\texttt{SmolLM3-3B-checkpoints} & 3B & 26 & 1246 & 38 \\
\texttt{apertus3-3b-64-nodes} & 3B & 13 & 1006 & 30 \\
\texttt{tiny-aya-base} & 3B & 1 & 1006 & 30 \\
\texttt{tiny-aya-earth} & 3B & 1 & 1006 & 30 \\
\texttt{tiny-aya-fire} & 3B & 1 & 1006 & 30 \\
\texttt{tiny-aya-global} & 3B & 1 & 1006 & 30 \\
\texttt{tiny-aya-water} & 3B & 1 & 1006 & 30 \\
\texttt{Qwen3-4B} & 4B & 1 & 1026 & 30 \\
\texttt{Qwen3-4B-Base} & 4B & 1 & 1006 & 30 \\
\texttt{gemma-3-4b-it} & 4B & 1 & 1006 & 30 \\
\texttt{gemma-3-4b-pt} & 4B & 1 & 1006 & 30 \\
\texttt{OLMo-2-1124-7B} & 7B & 1 & 1006 & 30 \\
\texttt{OLMo-2-1124-7B-DPO} & 7B & 1 & 1006 & 30 \\
\texttt{OLMo-2-1124-7B-Instruct} & 7B & 1 & 1006 & 30 \\
\texttt{OLMo-2-1124-7B-SFT} & 7B & 1 & 1006 & 30 \\
\texttt{Olmo-3-1025-7B} & 7B & 16 & 1070 & 34 \\
\texttt{Olmo-3-7B-Instruct} & 7B & 1 & 1021 & 30 \\
\texttt{Olmo-3-7B-Instruct-DPO} & 7B & 1 & 1006 & 30 \\
\texttt{Olmo-3-7B-Instruct-SFT} & 7B & 1 & 1006 & 30 \\
\texttt{Olmo-3-7B-Think} & 7B & 1 & 1006 & 30 \\
\texttt{Olmo-3-7B-Think-DPO} & 7B & 1 & 1006 & 30 \\
\texttt{Olmo-3-7B-Think-SFT} & 7B & 1 & 1006 & 30 \\
\texttt{Apertus-8B-2509} & 8B & 14 & 1070 & 34 \\
\texttt{Apertus-8B-Instruct-2509} & 8B & 1 & 1021 & 30 \\
\texttt{Qwen3-8B} & 8B & 1 & 1026 & 30 \\
\texttt{Qwen3-8B-Base} & 8B & 1 & 1006 & 30 \\
\texttt{aya-expanse-8b} & 8B & 1 & 1026 & 30 \\
\texttt{gemma-3-12b-it} & 12B & 1 & 1006 & 30 \\
\texttt{gemma-3-12b-pt} & 12B & 1 & 1006 & 30 \\
\texttt{OLMo-2-1124-13B} & 13B & 1 & 1006 & 30 \\
\texttt{OLMo-2-1124-13B-DPO} & 13B & 1 & 1006 & 30 \\
\texttt{OLMo-2-1124-13B-Instruct} & 13B & 1 & 1006 & 30 \\
\texttt{OLMo-2-1124-13B-SFT} & 13B & 1 & 1006 & 30 \\
\texttt{Qwen3-14B} & 14B & 1 & 1026 & 30 \\
\texttt{Qwen3-14B-Base} & 14B & 1 & 1006 & 30 \\
\texttt{gemma-3-27b-it} & 27B & 1 & 1006 & 30 \\
\texttt{gemma-3-27b-pt} & 27B & 1 & 1006 & 30 \\
\texttt{Qwen3-30B-A3B} & 30B-A3B & 1 & 1026 & 30 \\
\texttt{OLMo-2-0325-32B} & 32B & 1 & 1006 & 30 \\
\texttt{OLMo-2-0325-32B-DPO} & 32B & 1 & 1006 & 30 \\
\texttt{OLMo-2-0325-32B-Instruct} & 32B & 1 & 1006 & 30 \\
\texttt{OLMo-2-0325-32B-SFT} & 32B & 1 & 1006 & 30 \\
\texttt{Olmo-3-1125-32B} & 32B & 1 & 1006 & 30 \\
\texttt{Olmo-3-32B-Think} & 32B & 1 & 1006 & 30 \\
\texttt{Olmo-3-32B-Think-DPO} & 32B & 1 & 1006 & 30 \\
\texttt{Olmo-3-32B-Think-SFT} & 32B & 1 & 1006 & 30 \\
\texttt{Olmo-3.1-32B-Instruct} & 32B & 1 & 1006 & 30 \\
\texttt{Olmo-3.1-32B-Instruct-DPO} & 32B & 1 & 1006 & 30 \\
\texttt{Olmo-3.1-32B-Instruct-SFT} & 32B & 1 & 1006 & 30 \\
\texttt{aya-expanse-32b} & 32B & 1 & 1026 & 30 \\
\texttt{Apertus-70B-2509} & 70B & 1 & 1006 & 30 \\
\texttt{Qwen3-0.6B} & 600M & 1 & 1026 & 30 \\
\texttt{Qwen3-0.6B-Base} & 600M & 1 & 1006 & 30 \\
\texttt{apertus-0.6b-from8b-TOP256-long} & 600M & 10 & 1006 & 30 \\
\bottomrule
\end{tabular}
\end{table*}

\section{Signal and Noise}
\label{subsec:methodology-snr}

The signal-to-noise ratio (SNR) summarizes a benchmark's reliability as the ratio between signal, how well a benchmark separates models of similar compute, and noise, the variability in scores across training checkpoints or stochastic runs.
We define it as
\begin{equation}
\text{SNR} = \frac{\text{Signal}}{\text{Noise}}.
\end{equation}
A higher SNR means score differences across models are large relative to the benchmark's intrinsic variability. The framework is not tied to a particular signal or noise definition, so we evaluate several candidates and select the pair with the strongest empirical correlation to DA.

\textbf{Signal candidates.} Given the final checkpoints of training runs at similar compute $C_\text{final}$, we consider five candidate signal definitions to capture a different notion of how spread out model scores are. Variance is the average squared distance from the mean. Mean distance is the average pairwise distance between points. Relative standard deviation is the standard deviation divided by the mean. Star discrepancy is the largest difference between any point and the uniform distribution. Dispersion is the largest pairwise distance,
\begin{equation}
\text{Dispersion}(C_\text{final}) = \max_{i \neq j} \|c_i - c_j\|.
\end{equation}

\textbf{Noise candidates.} We consider four sources of noise that the framework can in principle account for. Seed noise is the standard deviation of the final checkpoint across training runs with different random seeds. Data order noise is the equivalent across runs that vary the order of training documents. Total variation is the average change in metric score across training checkpoints minus an improvement term. Checkpoint-to-checkpoint noise is proposed by \citet{heineman_signal_2025} as a proxy when the first three are too expensive to estimate at large scale and depends only on the final $n$ training checkpoints:
\begin{equation}
\text{Checkpoint noise} = \sigma\big(\{U(t_j)\}_{j=T-k+1}^{T}\big).
\end{equation}
On our ladder both checkpoint noise, over the last five evaluated checkpoints of each run, and seed noise, over the cells trained with three seeds, are measurable, and \S\ref{subsec:rq4} reports every intervention's effect against the seed noise.
We also consider a fifth noise definition, computable from a single evaluation run on any model. Benchmark noise is the average relative standard deviation of scores across $k = 5$ disjoint folds of the evaluation set:
\begin{equation}
\text{Noise}(M) = \frac{1}{|M|} \sum_{m \in M} \frac{\sqrt{\tfrac{1}{k}\sum_{i=1}^{k}(m_i - \bar{m})^2}}{\bar{m}}
\end{equation}
where $m_1, \dots, m_k$ are the per-fold scores of model $m$. The fold scores are obtained by partitioning the stored per-question outputs of a single evaluation run, so no inference is repeated and no intermediate checkpoints are needed.


\subsection{Measures of signal}
\label{app:measuring-signal}

When designing an measure of signal, we want to incorporate the uniformity of benchmark scores and the overall range of scores. Given the final checkpoints of training runs under similar compute spend $C_\text{final}$, we evaluate multiple approaches to measuring signal:

\begin{itemize}[itemsep=0pt,topsep=0pt]
    \item \textbf{Variance} measures average squared distance from the mean: $\text{Var}(C_\text{final}) = \frac{1}{n} \sum_{i=1}^{n} \|c_i - \bar{c}\|^2$
    \item \textbf{Mean distance} measures average pairwise distance between points: $\text{Mean Dist}(C_\text{final}) = \frac{2}{n(n-1)} \sum_{i < j} \|c_i - c_j\|$
    \item \textbf{Relative standard deviation}, or the coefficient of variation, measures the standard deviation divided by the mean: $\text{Rel. Std.}(C_\text{final}) = \frac{\sqrt{\text{Var}(C_\text{final})}}{\text{Mean}(C_\text{final})}$
    \item \textbf{Star Discrepancy} measures the largest difference between any point and the uniform distribution: $\text{Discrepancy}(C_\text{final}) = \sup_{t \in [0,1]} \left| \frac{1}{n} \sum_{i=1}^n \mathbf{1}\{c_i \leq t\} - t \right|$. 
    \item \textbf{Dispersion} measures the largest difference between any two points, or the largest unfilled space in the range of performance: $\text{Dispersion}(C_\text{final}) = \max_{i \neq j} \|c_i - c_j\|$. 
\end{itemize}


\subsection{Measures of Modeling noise{}}
\label{app:measuring-noise}

\textbf{Seed noise{}.} To measure the noise{} introduced from changing the random seed initialization between training runs, we can compute the standard deviation of the final checkpoint from multiple training runs with different random seeds. To estimate seed noise{}, we train $M$ models using the same configuration, and average the scores over the final $n$ checkpoints of $T$ total training checkpoints to smooth the checkpoint-to-checkpoint noise{}, then compute the standard deviation:
\begin{equation}
 \text{Seed noise{}}(M) = \sigma(M),\quad M_i = \textstyle\frac{1}{n} \sum_{j=T-n+1}^{T} U(t_j)
\end{equation}

\textbf{Data Order noise{}.} This is noise{} introduced from changing the order of sampled documents from the training data. We estimate the data order noise{} using the same method as seed noise{}.

\textbf{Total Variation.} To measure the checkpoint-to-checkpoint noise{} throughout an entire training run, we measure the total variation of the intermediate training checkpoints on the downstream benchmark. We measure total variation as the average change in metric score across $T$ training checkpoints minus an improvement term:
\label{app:total-variation}
\begin{equation}
\text{Total Variation} = \textstyle\frac{1}{T} \sum_{t=1}^{T} |U(t) - U(t-1)| - \frac{1}{T}(U(T) - U(0))
\end{equation}

\textbf{Checkpoint-to-checkpoint noise{}.} 
Calculating the above sources of noise{} are either too expensive to estimate at large scales (e.g., training LLMs by varying the random seed) or difficult to run (e.g., evaluating every checkpoint on an LLM training curve). Instead, we propose an estimate measuring only the noise{} of the final $n$ training checkpoints of training:
\begin{equation}
\text{Checkpoint-to-checkpoint noise{}} = \textstyle\sigma\left(\{U(t_j)\}_{j=T-k+1}^{T}\right)
\end{equation}

\textbf{Benchmark noise.} We propose to measure noise as the relative standard deviation of scores across $k = 5$ disjoint folds of the evaluation set, averaged across models:

$$\text{Noise}(M) = \frac{1}{|M|} \sum_{m \in M} \frac{\sqrt{\frac{1}{k}\sum_{i=1}^{k}(m_i - \bar{m})^2}}{\bar{m}}$$

where $m_1, \dots, m_k$ are the per-fold scores of model $m$. We adopt the benchmark noise over the original checkpoint-to-checkpoint noise for two reasons. First, it is computable from a single evaluation run, since fold scores are derived by partitioning stored per-question outputs without rerunning inference. This makes the framework applicable to any public model, including those without released intermediate checkpoints. Second, it correlates more strongly with decision accuracy than checkpoint noise ($R = 0.808$ vs. $R = 0.760$).


\end{document}
